% Generated draft: requires mathematical and source review.
\documentclass[11pt]{article}
\usepackage[margin=25mm]{geometry}
\usepackage{fontspec}
\setmainfont{Latin Modern Roman}
\usepackage{amsmath,amssymb,mathtools,unicode-math}
\setmathfont{Latin Modern Math}
\usepackage{enumitem,needspace,xurl,hyperref}
\setlength{\parindent}{0pt}
\setlength{\parskip}{5pt}
\setlength{\emergencystretch}{3em}
\newcommand{\N}{\mathbb{N}}
\newcommand{\Z}{\mathbb{Z}}
\newcommand{\Q}{\mathbb{Q}}
\newcommand{\R}{\mathbb{R}}
\newcommand{\C}{\mathbb{C}}
\newcommand{\F}{\mathbb{F}}
\newcommand{\A}{\mathbb{A}}
\newcommand{\PP}{\mathbb{P}}
\newcommand{\E}{\mathbb{E}}
\newcommand{\eps}{\varepsilon}
\newcommand{\dd}{\,\mathrm{d}}
\begin{document}
\section*{UnsolvedMath ID: 800953}
\textbf{Source problem code:} catalogue entry
\subsection{Definitions and mathematical statement}
A closed manifold means a compact smooth manifold without boundary. For a Riemannian metric \(g\) on \(M\), let \(K_g(\sigma)\) denote the sectional curvature of a two-dimensional linear subspace \(\sigma\subset T_xM\). Write
\[
SM=\{(x,v)\in TM:g_x(v,v)=1\}
\]
for the unit tangent bundle. The geodesic flow \(\varphi_t:SM\to SM\) is defined by
\[
\varphi_t(x,v)=(\gamma_{x,v}(t),\dot\gamma_{x,v}(t)),
\]
where \(\gamma_{x,v}\) is the geodesic with initial data \((x,v)\). With respect to the normalized Liouville measure \(\mu_L\), ergodicity means that every measurable set \(B\subseteq SM\) satisfying \(\mu_L(\varphi_t^{-1}(B)\mathbin{\triangle}B)=0\) for every \(t\in\R\) has \(\mu_L(B)\in\{0,1\}\).

Under the conventional interpretation of ergodicity with respect to Liouville measure and of positive sectional curvature in dimension at least two, the question is: does there exist a closed Riemannian manifold \((M,g)\), of some dimension \(n\geq2\), such that
\[
K_g(\sigma)>0
\quad\text{for every }x\in M\text{ and every two-dimensional }\sigma\subset T_xM,
\]
and whose geodesic flow on \(SM\) is ergodic with respect to \(\mu_L\)?

The supplied source does not explicitly specify the invariant measure or the admissible dimensions; the conventions used above require confirmation against the intended source formulation. The source describes the general problem as still wide open despite reports of its solution on page 87 of [Berger2000] and in Section 10.9 of [Berger2003]. This is a supplied historical assessment, not an independently verified conclusion.
\subsection{Short English statement}
Does some closed manifold admit a Riemannian metric with strictly positive sectional curvature everywhere and an ergodic geodesic flow? The draft uses the conventional Liouville measure and dimension at least two; these conventions are not explicit in the supplied source.
\subsection{Sources}
\begin{itemize}
\item Supplied local UnsolvedMath registry record, permanent ID \(800953\), exact source problem code \texttt{catalogue entry}; title: \emph{Open Problems and Questions about Geodesics --- Question 33}. The registry attributes the problem to Burns and Matveev, gives the proposed year as 2019, and labels it \texttt{open}. That label is supplied metadata, not independently verified mathematical status.
\item Burns and Matveev, \emph{Open Problems and Questions about Geodesics} (2019, as identified in the supplied record), Question 33, \url{https://arxiv.org/abs/1308.5417}. The supplied provenance records access on 2026-10-16 and extraction by \texttt{arxiv\_latex}; no independent source visit or verification is asserted here.
\item [Berger2000], page 87, and [Berger2003], Section 10.9: cited in the supplied statement as places reporting a solution. Full bibliographic details and the content of those reports were not supplied. The supplied statement nevertheless describes the general problem as still wide open.
\end{itemize}
\end{document}
